
\documentclass[a4paper,12pt]{article}
\usepackage[margin=2cm]{geometry}

% Pacchetti utili di base
\usepackage[utf8]{inputenc}   % codifica dei caratteri
\usepackage[T1]{fontenc}      % font encoding
\usepackage[italian]{babel}   % lingua italiana
\usepackage{graphicx}         % immagini
\usepackage{hyperref}         % link cliccabili
\usepackage{amssymb}
\usepackage{float}
\usepackage{amsmath}
\usepackage{enumitem}
\usepackage{bm}               % vettori e matrici in grassetto
\usepackage{array}            % tabelle
\usepackage{xcolor}
\usepackage{tikz}             % disegni e alberi
\usetikzlibrary{positioning,arrows.meta,shapes.geometric}
\usepackage{pgfplots}
\pgfplotsset{compat=1.18}

\graphicspath{{img/}}

% Liste piu' compatte e allineate al testo
\setlist[itemize]{topsep=0.4em, itemsep=0.25em, parsep=0pt, leftmargin=1.6em}
\setlist[enumerate]{topsep=0.4em, itemsep=0.25em, parsep=0pt, leftmargin=1.9em}

% Notazione ricorrente
\newcommand{\R}{\mathbb{R}}
\newcommand{\E}{\mathbb{E}}
\newcommand{\Prob}{\mathbb{P}}
\newcommand{\x}{\bm{x}}
\newcommand{\w}{\bm{w}}
\newcommand{\X}{\bm{X}}
\newcommand{\yhat}{\hat{y}}
\newcommand{\norm}[1]{\left\lVert #1 \right\rVert}
\newcommand{\argmin}{\operatorname*{arg\,min}}
\newcommand{\argmax}{\operatorname*{arg\,max}}

% --- Presentazione uniforme delle figure -------------------------------------
% Stessa aria sopra/sotto per ogni immagine o disegno, centratura che funziona
% anche dentro le liste, nessuno spostamento dovuto ai float.
\newlength{\figgap}
\setlength{\figgap}{1.1em}

\newcommand{\figura}[1]{%
  \par\addvspace{\figgap}%
  {\centering\noindent#1\par}%
  \addvspace{\figgap}%
  \noindent\ignorespaces
}
\newcommand{\imgfig}[2][0.55]{\figura{\includegraphics[width=#1\linewidth]{#2}}}

% Voce di elenco con testo a sinistra e immagine a destra:
% #1 = larghezza relativa dell'immagine, #2 = testo, #3 = file immagine
\newcommand{\itemimg}[3][0.85]{%
  \begin{minipage}[t]{0.70\linewidth}%
    #2%
  \end{minipage}%
  \hfill
  \begin{minipage}[t]{0.26\linewidth}%
    \vspace{0pt}\centering\includegraphics[width=#1\linewidth]{#3}%
  \end{minipage}%
}

% Evidenziatura usata nelle slide
\newcommand{\hl}[1]{\textcolor{orange}{#1}}

% --- Albero di decisione di esempio (Play Tennis) ----------------------------
\newcommand{\alberoTennis}{%
\begin{tikzpicture}[
    level 1/.style={sibling distance=42mm, level distance=15mm},
    level 2/.style={sibling distance=22mm, level distance=15mm},
    attr/.style={draw, rounded corners=2pt, text=blue, font=\footnotesize,
                 inner sep=2pt},
    si/.style={text=green!50!black, font=\footnotesize},
    no/.style={text=red, font=\footnotesize},
    el/.style={font=\scriptsize, midway, fill=white, inner sep=1pt}
  ]
  \node[attr] {Outlook}
    child { node[attr] {Humidity}
      child { node[no] {No}   edge from parent node[el,left]  {High}   }
      child { node[si] {Yes}  edge from parent node[el,right] {Normal} }
      edge from parent node[el,left] {Sunny} }
    child { node[si] {Yes} edge from parent node[el,left,xshift=-2mm] {Overcast} }
    child { node[attr] {Wind}
      child { node[no] {No}  edge from parent node[el,left]  {Strong} }
      child { node[si] {Yes} edge from parent node[el,right] {Weak}   }
      edge from parent node[el,right] {Rain} };
\end{tikzpicture}%
}

% --- I due alberi candidati A e B --------------------------------------------
\newcommand{\alberiAB}{%
\begin{tikzpicture}[
    nd/.style={circle, draw, minimum size=7mm, inner sep=0pt, font=\small},
    cnt/.style={font=\scriptsize},
    el/.style={font=\scriptsize, midway, fill=white, inner sep=1pt}
  ]
  % albero con radice A
  \node[nd] (a)  at (0,0)       {A};
  \node[cnt, left=2mm of a] {$[29+,35-]$};
  \node[nd] (al) at (-1.3,-1.7) {};
  \node[nd] (ar) at ( 1.3,-1.7) {};
  \node[cnt, left=2mm of al] {$[21+,5-]$};
  \node[cnt, left=2mm of ar] {$[8+,30-]$};
  \draw (a) -- (al) node[el, left]  {T};
  \draw (a) -- (ar) node[el, right] {F};

  % albero con radice B
  \node[nd] (b)  at (7,0)       {B};
  \node[cnt, left=2mm of b] {$[29+,35-]$};
  \node[nd] (bl) at (5.7,-1.7) {};
  \node[nd] (br) at (8.3,-1.7) {};
  \node[cnt, left=2mm of bl] {$[18+,33-]$};
  \node[cnt, left=2mm of br] {$[11+,2-]$};
  \draw (b) -- (bl) node[el, left]  {T};
  \draw (b) -- (br) node[el, right] {F};
\end{tikzpicture}%
}

% Informazioni sul documento
\title{Introduzione all'apprendimento automatico}
\author{Nicolò Giuliani}
\date{}


\begin{document}

\maketitle
\tableofcontents
\newpage

\setcounter{section}{-1}   % la prima \section è numerata 0

\section{Introduzione}

\subsection{Di cosa si occupa il Machine Learning?}
Si occupa di risolvere problemi complessi tramite tecniche algoritmiche specifiche. L'approccio generale del ML è
\begin{itemize}
    \item definire una classe di modelli del problema su insieme di \textbf{parametri} $\Theta$
    \item definire una metrica di valutazione, ovvero una misura dell'\textbf{errore} dei modelli
    \item modificare i parametri $\Theta$ per minimizzare l'errore sui \textbf{dati di training}
\end{itemize}
Quindi possiamo dire che \texttt{il machine learning è un procedimento di ottimizzazione}.
\imgfig[0.3]{2026-09-24_11-35-45.png}
Il ML tratta di problemi di \textbf{ottimizzazione}!
\begin{itemize}
    \item Il \texttt{training set} fa sì che il tuning dei parametri è dato dalla esperienza passata, quindi una sorta di \textbf{apprendimento}.
    \item La \texttt{soluzione} spesso non esiste una soluzione in forma chiusa
\end{itemize}
Utilizziamo \textbf{tecniche iterative} per approssimare iterativamente la soluzione.

\subsubsection{Predizioni}
Attraverso modelli di ML possiamo anche stabilire correlazioni, clusters, $\dots$
\imgfig[0.55]{2026-09-24_11-45-42.png}

\subsubsection{Classi di problemi di apprendimento}
\begin{itemize}
    \item \itemimg{%
        Approfondimento supervisionato: inputs + outputs (etichette)
        \begin{itemize}[topsep=0.2em, partopsep=0pt]
            \item classificazione
            \item regressione
        \end{itemize}}{2026-09-24_11-49-31.png}

    \item \itemimg{%
        Approfondimento non supervisionato: inputs
        \begin{itemize}[topsep=0.2em, partopsep=0pt]
            \item clustering
            \item analisi delle componenti
            \item autoencoding
        \end{itemize}}{2026-09-24_11-49-58.png}

    \item \itemimg{%
        Approfondimento con rinforzo
        \begin{itemize}[topsep=0.2em, partopsep=0pt]
            \item apprendere comportamenti
            \item ``model free'' planning
        \end{itemize}}{2026-09-24_11-50-22.png}
\end{itemize}

\subsubsection*{Molte tecniche differenti}
\begin{itemize}
    \item Modi diversi per definire i modelli
    \begin{itemize}
        \item alberi di decisione
        \item modelli lineari
        \item reti neurali
    \end{itemize}
    \item Differenti funzioni di loss
    \begin{itemize}
        \item errore quadratico
        \item logistic loss
        \item cross entropy
        \item similarità del coseno
        \item margine massimo
    \end{itemize}
\end{itemize}
\imgfig[0.50]{2026-09-24_11-53-02.png}

\subsection{Features}
Ogni informazione realativa a un dato che ne descrive una sua proprietà è feature (caratteristica).
\begin{itemize}
    \item Le features sono gli input del processo di apprendimento.
    \item L'apprendimento è molto sensibile alla scelta delle features.
\end{itemize}
Determinare delle buone caratteristiche è complesso.
\begin{itemize}
    \item \textbf{Approccio tradizionale}: determinare durante il preprocessing delle buone features, e applicare un metodo semplice ma robusto di apprendimento, ad esempio qualche tecnica lineare.
    \item \textbf{Approccio ``deep''}: fornire in input dati grezzi e delegare alla macchina il compito di sintetizzare nuove features (rappresentatione interna).
\end{itemize}
\texttt{Deep learning} è implementato attraverso reti neurali profonde.\\
La rete neurale è profonda quando si hanno layers multipli:
\imgfig[0.55]{2026-09-24_11-59-11.png}
Ogni layer sintetizza nuove features in funzione delle precedenti.

\subsubsection*{Relazione tra le aree di ricerca}
\imgfig[0.55]{2026-09-24_12-00-05.png}
L'AI, machine learning, deep learning:
\begin{itemize}
    \item \textbf{Sistemi basati su conoscenza}: prendete un esperto, chiedetegli come risolverebbe il problema e poi provate a mimare il suo comportamento mediante regole logiche
    \item \textbf{Machine-Learning tradizionale}: prendete un esperto, chiedetegli di quali features avrebbe bisogno per risolvere un dato problema, raccogliete un insieme di osservazioni e lasciate alla macchina il compito di apprendere l'associazione
    \item \textbf{Deep-Learning}: liberatevi dell'esperto
\end{itemize}

\subsubsection*{Componenti addestrati mediante apprendimento}
\imgfig[0.55]{2026-09-24_12-01-52.png}

\subsection{Alberi di decisione}
Definiamo \textbf{traning set} un insieme di \texttt{esempi di allenamento}
\begin{align*}
    \langle x^{(i)}, y^{(i)}\rangle
\end{align*}
con
\begin{itemize}
    \item $x^{(i)} \in X$ insieme degli inputs
    \item $y^{(i)} \in Y$ insieme delgi outputs, \textbf{groudn truth}
    \item $i$ indice dell'istanza di training
\end{itemize}
Definiamo \textbf{problema} trovare la funzione che ``mappa'' $x^{(i)}$ a $y^{(i)}$
\begin{itemize}
    \item $Y$ discreto: problema di classificazione (predizione di una classe)
    \item $Y$ continuo: problema di regressione (predizione di un valore)
\end{itemize}
Le tecniche di Machine Learning richiedono la scelta preliminare di uno \texttt{spazio di funzioni} $H$, all'interno del quale cercare la soluzione che \texttt{meglio approssima} i dati di training. \\
\underline{Un modello è un modo di specificare e calcolare una funzione nello}\\
\underline{spazio delle ipotesi $H$.}

\subsubsection{Alberi di decisione: un esempio}
È un buon giorno per giocare a tennis?
\begin{align*}
    H : \text{Outlook} \times \text{Humidity} \times \text{Wind} \times \text{Temp} \to \text{Play Tennis?}
\end{align*}
\figura{\alberoTennis}
Ogni nodo testa un attributo (feature) $X$\\
Ogni arco uscente da un nodo corrisponde a uno dei possibili valori \hl{discreti} di $X$\\
Ogni foglia predice la riposta $Y$ (o un probabilità $P(Y|X)$)
\begin{itemize}
    \item Insieme di input $X$:\\
    ogni istanza $x \in X$ è un vettore di features del seguente tipo:\\
    $\langle\, \text{Humidity=high},\ \text{Wind=weak},\ \text{Outlook=rain},\ \text{Temp=hot} \,\rangle$
    \item Funzione target $f : X \to Y$\\
    $Y$ assume valori discreti (booleani)
    \item Spazio delle ipotesi $H = \{h \mid h : X \to Y\}$ (nessuna restrizione)
    \begin{itemize}
        \item vogliamo modellare $h \in H$ con un albero di decisione
        \item ogni istanza $x \in X$ definisce un cammino nell'albero che conduce a una foglia etichettata con $y \in Y$, corrispondente alla soluzione predetta
    \end{itemize}
\end{itemize}
\figura{%
  \begin{tabular}{|l|l|l|l|l|}
    \hline
    Outlook  & Temp & Humidity & Wind   & Play \\
    \hline
    Sunny    & Hot  & High     & Weak   & No  \\
    \hline
    Sunny    & Hot  & High     & Strong & No  \\
    \hline
    Overcast & Hot  & High     & Weak   & Yes \\
    \hline
    Rain     & Mild & High     & Weak   & Yes \\
    \hline
    Rain     & Cool & Normal   & Weak   & Yes \\
    \hline
    Rain     & Cool & Normal   & Strong & No  \\
    \hline
    Overcast & Cool & Normal   & Strong & Yes \\
    \hline
    Sunny    & Mild & High     & Weak   & No  \\
    \hline
    Sunny    & Cool & Normal   & Weak   & Yes \\
    \hline
    Rain     & Mild & Normal   & Weak   & Yes \\
    \hline
    Sunny    & Mild & Normal   & Strong & Yes \\
    \hline
    Overcast & Mild & High     & Strong & Yes \\
    \hline
    Overcast & Hot  & Normal   & Weak   & Yes \\
    \hline
    Rain     & Mild & High     & Strong & No  \\
    \hline
  \end{tabular}%
}

\paragraph{Expressività del modello}
Definiamo $X_1 \times X_2 \times \dots \times X_n$ con $X_i = \{\text{true, false}\}$ \\
Possiamo esprimere $Y=X_2 \lor X_5$?

\subsubsection{Costruzione Top-down induttiva}
Main loop:
\begin{enumerate}
    \item assegnare al nodo corrente il ``miglior'' attributo $X_i$
    \item creare un nodo figlio per ogni possibile valore di $X_i$ e propagare i dati verso i figli a seconda del loro valore
    \item per ogni nodo figlio, se tutti i dati del training set associati al nodo hanno una stessa etichetta $y$, marcare il nodo come foglia con etichetta $y$, altrimenti ripetere dal primo punto
\end{enumerate}
Ma quale è il ``miglior'' attributo?
\figura{\alberiAB}
Definiamo \textbf{entropia} di una variabile aleatoria $X$
\begin{align*}
    H(X)=- \sum_{i=1}^{n} P(X=i)\log_2 P(X=i)
\end{align*}
dove $n$ è il numero dei possibili valori di $n$.

\figura{%
  \begin{minipage}[c]{0.48\linewidth}
    L'entropia misura il \hl{grado di impurità} (disordine) della informazione.
    È massima ($\log n$) quando $X$ è uniformemente distribuita tra tutti i
    suoi $n$ valori, e minima ($0$) quando è concentrata su un singolo valore.
  \end{minipage}%
  \hfill
  \begin{minipage}[c]{0.46\linewidth}
    \centering
    \begin{tikzpicture}
      \begin{axis}[
          width=\linewidth, height=0.72\linewidth,
          xlabel={$\Pr(X=1)$}, ylabel={$H(X)$},
          xmin=0, xmax=1, ymin=0, ymax=1,
          xtick={0,0.5,1}, ytick={0,0.5,1},
          grid=both, grid style={gray!25, very thin},
          tick label style={font=\scriptsize},
          label style={font=\scriptsize},
          axis lines=box,
        ]
        \addplot[blue, thick, samples=200, domain=0.0005:0.9995]
          {-x*log2(x)-(1-x)*log2(1-x)};
      \end{axis}
    \end{tikzpicture}
  \end{minipage}%
}

\subsubsection{Teoria della Informazione (Shannon)}
L'Entropia è la quantità media di informazione prodotta da una sorgente stocastica di dati.
\begin{itemize}
    \item Un evento con entropia 1 non trasmette Informazioni
    \begin{align*}
        I(1)=0
    \end{align*}
    \item Dati due evento indipendenti con probabilità $p_1, p_2$ la probabilità congiunta è $p_1p_2$ ma l'informazione acquisita è la somma delle informazioni dei due eventi indipendenti
    \begin{align*}
        I(p_1p_2)= I(p_1)+I(p_2)
    \end{align*}
    \item Ci aspettiamo che l'informazione sia antimonotona rispetto alla probabilità ed è quindi naturale definire
    \begin{align*}
        I(p)=-\log(p)
    \end{align*}
\end{itemize}

\subsubsection*{Teoria dei codici (Shannon)}
L'Entropia misura il numero medio di bits richiesti per trasmettere il valore prodotto da una sorgente stocastica $X$. \\
\underline{Se gli eventi non hanno la stessa probabilità possiamo migliorare la codifica!!}

\paragraph{Entropia di $X$}
\begin{align*}
    H(X) = -\sum_{i=1}^{n} P(X=i)\log_2 P(X=i)
\end{align*}

\paragraph{Entropia Condizionale di $X$ dato uno specifico $Y=v$}
\begin{align*}
    H(X|Y=v) = -\sum_{i=1}^{n} P(X=i|Y=v)\log_2 P(X=i|Y=v)
\end{align*}

\paragraph{Entropia Condizionale di $X$ dato $Y$}
(media ponderata sugli $m$ possibili valori di $Y$)
\begin{align*}
    H(X|Y) = \sum_{v=1}^{m} P(Y=v)H(X|Y=v)
\end{align*}

\paragraph{Guadagno Informativo tra $X$ e $Y$}
\begin{align*}
    I(X,Y) = H(X)-H(X|Y) = H(Y)-H(Y|X)
\end{align*}

\subsubsection{Meglio A o B?}
Misuriamo la riduzione di entropia della variabile target $Y$ conseguente alla conoscenza di un qualche attributo $X$, cioè il guadagno informativo $I(Y,X)$ tra $Y$ e $X$.
\figura{\alberiAB}
\begin{align*}
    H(Y)       &= -(29/64)\cdot\log_2(29/64)-(35/64)\cdot\log_2(35/64) = .994\\
    H(Y|A=T)   &= -(21/26)\cdot\log_2(21/26)-(5/26)\cdot\log_2(5/26)   = .706\\
    H(Y|A=F)   &= -(8/38)\cdot\log_2(8/38)-(30/38)\cdot\log_2(30/38)   = .742\\
    H(Y|A)     &= .706\cdot 26/64 + .742\cdot 38/64 = .726\\
    I(Y,A)     &= H(Y)-H(Y|A) = .994-.726 = .288
\end{align*}
Per $B$ otteniamo $H(Y|B) = .872$ e $I(Y,B) = .122$
\begin{center}
    \textcolor{red}{Quindi, A è migliore!}
\end{center}

\subsubsection{Il caso continuo}
Quando gli attibuti sono continui, prendiamo le decisioni in base a opportune \hl{soglie} (thresholds):
\figura{%
  \begin{tikzpicture}[
      nd/.style={ellipse, draw, minimum width=16mm, minimum height=11mm,
                 inner sep=1pt, font=\small},
      el/.style={font=\small}
    ]
    \node[nd] (t) at (0,0) {$A<\Theta$};
    \draw[-{Stealth[length=2.5mm]}] (t) -- (-1.6,-1.6)
      node[el, text=green!50!black, pos=0.15, left] {T};
    \draw[-{Stealth[length=2.5mm]}] (t) -- ( 1.6,-1.6)
      node[el, text=red, pos=0.15, right] {F};
  \end{tikzpicture}%
}
\begin{itemize}
    \item confrontiamo le varie soglie in base al loro guadagno informativo
    \item ma come scegliere le soglie candidate?
    \begin{itemize}
        \item sampling a intervalli discreti prefissati
        \item ordinare il train set rispetto a un dato attributo e scegliere come soglie i valori medi tra dati consecutivi
    \end{itemize}
\end{itemize}

\subsection{Overfitting}
\end{document}
